% Volume VIII: Applications and Case Structures
% Part XVI. Failure Modes of the Theory
% Chapter 91: When Procedure Protects Power
% Spine: proof-spines/ch091-spine.md

\chapter{When Procedure Protects Power}
\label{ch:ch091}

Due process is not equally available to all parties. Wealthy actors can convert evidentiary rigor into indefinite delay, procedural complexity, and reputation management, while poorer actors often experience rapid coercion under the same nominal system. This chapter addresses that asymmetry directly so that the book's defense of disciplined suspicion does not collapse into a romance of actually existing legal and administrative institutions.

The focus is on the failure mode in which procedure works exactly as designed for the protection of the already powerful. Rules about pleading, counsel, delay, discovery, confidentiality, and burden can be formally neutral while remaining usable only by those who can afford time, expertise, and retaliation risk. The chapter therefore asks when procedural friction serves truth and when it merely converts hierarchy into the appearance of fairness.
